\begin{textsample}{POS Dim 1 – mg – Score 51.00 – mg/mg\_em\_13.txt}  \label{ex:f1_pos_233}
1.5.1 Desenvolvimento das \textbf{competências} gerais na etapa do Ensino Médio As Dez \textbf{Competências} Gerais a serem desenvolvidas pelos estudantes ao longo da Educação Básica apresentam uma \textbf{inter-relação} e desdobram -se na ação pedagógica, “ articulando -se na construção de conhecimentos, no desenvolvimento de habilidades e na formação de atitudes e valores ” ( BRASIL, 2018, p.8 ) essenciais para a vida no \textbf{século} XXI.
Nesse sentido, de forma articulada e progressiva, as propostas e práticas pedagógicas devem consubstanciar os direitos de aprendizagem e desenvolvimento entre as etapas da Educação Básica de forma que, no Ensino Médio, o estudante possa ampliar e aprofundar os conhecimentos com autonomia, responsabilidade, ética e consciência \textbf{socioambiental}.
A \textbf{Competência} 2 refere -se à criatividade e ao pensamento científico e \textbf{crítico} e, no âmbito pedagógico, as articulações e \textbf{progressões} devem propor exploração de ideias e suas conexões, criação de processos de investigação, formulação e desenvolvimento de hipóteses, \textbf{análise} de dados, lógica, \textbf{raciocínio}, avaliação de processos cognitivos, explicação de evidências e síntese. \textbf{Espera} -se que o estudante faça conexões adequadas e claras entre as ideias mais amplas, utilizando estratégias ou \textbf{caminhos} diferentes, dominando o processo de articulação entre as ideias mais \textbf{complexas}, resolvendo situações e propondo soluções criativas. \textbf{Espera} -se que ele elabore plano de investigação, abordando vários aspectos a serem analisados, formulando hipóteses, \textbf{verificando} -as com os dados analisados de forma científica, interpretando -os de forma coesa e clara, com a utilização de diferentes estratégias e conceitos analíticos, sendo capaz de avaliar e se posicionar de forma científica e ética, explicando as conclusões lógicas geradas a partir da investigação.
Na \textbf{Competência} 3, a \textbf{abordagem} pedagógica, com \textbf{foco} na \textbf{progressão} sistemática e articulada, pauta -se no repertório cultural, na identidade e diversidade cultural, propiciando ao estudante a fruição e participação em contextos artísticos e culturais, expressando sentimentos e experiências por meio das artes. \textbf{Espera} -se que ele possa desenvolver o senso de pertencimento a grupos sociais e culturais distintos, reconhecendo o significado de manifestações e \textbf{eventos} culturais e a \textbf{influência} e contribuições destes na construção da identidade individual e coletiva.
Neste sentido, ao longo dos três anos do Ensino Médio, objetiva -se que o estudante, por meio das artes, compreenda e valorize os contextos culturais, históricos, ambientais, aplicando conhecimentos e habilidades para expressar experiências, ideias e percepções.
Que analise obras criativas com mais \textbf{profundidade}, utilizando -as em outros contextos com finalidade e público diferentes, ampliando sua \textbf{compreensão} acerca de sua identidade e sua cultura e a de outrem, refletindo e agindo de forma flexível, respeitosa e acolhedora sobre questões relacionadas a valores e diferentes visões de mundo.
Desta forma, busca -se promover a convivência entre grupos distintos com valorização de intercâmbio \textbf{interculturais} e o respeito às diversas identidades, às manifestações, às culturas.
A \textbf{Competência} 4 \textbf{diz} respeito à comunicação, ao desenvolvimento da expressão, da escuta, da discussão e dos \textbf{multiletramentos}.
Nessa perspectiva, \textbf{espera} -se que o estudante, ao \textbf{final} do terceiro ano do Ensino Médio, compreenda e analise os discursos, expresse com clareza e coerência suas ideias, opiniões e sentimentos mais \textbf{complexos}, compartilhando informações e experiências com grupos de seu convívio ou outros.
Ademais, que domine os aspectos \textbf{discursivos} da comunicação verbal, com \textbf{vistas} a assegurar o entendimento mútuo, principalmente sobre temas mais amplos, preservando o \textbf{foco} do \textbf{debate}, utilizando estratégias de resumo e perguntas.
E que, também, examine \textbf{argumentos} que sustentem suas posições e os divergentes, construindo trajetos de discussões conectados com as ideias apresentadas.
E, ainda, que utilize plataformas analógicas e \textbf{digitais} mais \textbf{complexas} para se comunicar, fazendo uso de recursos, ferramentas e linguagens que apresentam uma diversidade de representações e/ou semioses, como áudio, textos, gráficos, linguagem \textbf{cartográfica}, matemática, artística, corporal ; dentre outras multimodalidades.
Já a \textbf{Competência} 5 pauta -se na cultura \textbf{digital}, pela qual se propõe o desenvolvimento do pensamento computacional, a computação, a programação, a cultura e o mundo \textbf{digital}, para produzir e aprender.
Nesse contexto, a partir de propostas e práticas pedagógicas \textbf{articuladas} e inter/transdisciplinares, \textbf{espera} -se que o estudante utilize recursos tecnológicos para projetar, desenvolver e apresentar produtos cada vez mais sofisticados, demonstrando conhecimentos na aplicação de dispositivos móveis, desenvolvendo e implementando \textbf{aplicativos} computacionais.
Também, que discuta os impactos da \textbf{tecnologia} nas comunicações e nas relações sociais e culturais, analisando, de forma \textbf{crítica} e ética, os comportamentos adequados ou não em redes sociais.
E que utilize a linguagem de programação para solucionar \textbf{problemas}, avaliando as desvantagens e potencialidades de diferentes algoritmos.
A \textbf{competência} 6 refere -se ao trabalho e projeto de vida.
Na primeira dimensão, o trabalho pedagógico deve considerar a \textbf{compreensão} sobre o mundo do trabalho e a preparação para ele, buscando desenvolver nos estudantes a criticidade e uma visão ampliada sobre desafios, relações e oportunidades associados ao mundo do trabalho na \textbf{contemporaneidade}, levando -os a reconhecer o valor do trabalho para a transformação social e a realização pessoal.
E, ainda, deve buscar que os estudantes desenvolvam a \textbf{capacidade} de analisar suas potencialidades, aspirações e desafios para escolhas mais assertivas, estabelecendo metas, inclusive com \textbf{projeções} financeiras, para a vida profissional presente e futura.
O capítulo 16, Texto \textbf{Introdutório} do Itinerário Formativo Educação Profissional e Técnica, aborda este tema mais detalhadamente.
Na segunda dimensão, as propostas e práticas docentes devem propiciar aos estudantes a oportunidade de desenvolver \textbf{capacidade} de lidar com desejos, \textbf{fragilidades} pessoais, frustrações e adversidades no desenvolvimento e alcance de seus objetivos e metas, \textbf{aprimorando} suas potencialidades, adquirindo autoconfiança, determinação e percebendo as adversidades e desafios como forma de crescimento e aperfeiçoamento das habilidades inter e intrapessoais, refletindo sobre seu desenvolvimento e suas metas e objetivos.
O capítulo 4, Projeto de Vida, traz uma \textbf{abordagem} mais aprofundada sobre essa dimensão.
A \textbf{Competência} 7 está relacionada à argumentação e à consciência global.
Busca -se desenvolver estratégias argumentativas, inferências e pontos de \textbf{vista}, bem como a consciência \textbf{socioambiental} numa perspectiva global.
Nesse sentido, é \textbf{imprescindível} a articulação entre os diferentes componentes curriculares de todas as áreas do conhecimento para um trabalho pedagógico sistemático e progressivo no alcance de tais premissas. \textbf{Espera} -se que os estudantes possam desenvolver sua argumentação por meio de estratégias argumentativas sólidas, expressando seu ponto de \textbf{vista} de forma clara, \textbf{ordenada}, coerente e compreensível, sendo capaz de confrontar pontos de \textbf{vistas} de forma respeitosa, com escuta a \textbf{posicionamentos} diferentes e aprendizagem a partir desse confronto.
Ademais, que eles desenvolvam a \textbf{capacidade} de inferir com pertinência, clareza e coerência.
Também, que os estudantes se interessem por questões globais, compreendam a \textbf{inter-relação} entre os desafios e \textbf{problemas} locais e mundiais, adquirindo uma visão ampliada sobre questões sociais, ambientais e do desenvolvimento da \textbf{ciência} e uma atitude respeitosa nas proposições de ações positivas e engajamento na promoção dos direitos humanos e da sustentabilidade ambiental.
A \textbf{Competência} 8 relaciona -se ao autoconhecimento e autocuidado, pautando -se no desenvolvimento da autoestima, da autoconsciência, da saúde física, do equilíbrio emocional, da reflexão e metarreflexão.
Para isso, por meio de projetos interdisciplinares, por exemplo, é essencial que as propostas de trabalho docente ofereçam ao estudante a oportunidade de vivenciar situações em que possa desenvolver -se, conhecendo a si mesmo em diversos contextos, buscando seu bem-estar e qualidade de vida para a manutenção da saúde física e do equilíbrio emocional em situações desafiadoras, inesperadas ou difíceis, compreendendo com mais \textbf{profundidade} suas emoções, reações e sentimentos em contextos distintos.
Também nesse viés, é importante que o estudante reflita sobre suas estratégias de escolha e sobre sua própria \textbf{maneira} de pensar, além de ser capaz de mudar o \textbf{foco} de atenção e responder a múltiplas \textbf{demandas} ou \textbf{tarefas} das \textbf{juventudes}.
Ou seja, ele será capaz de atender a diversas solicitações em contextos \textbf{complexos} que \textbf{exigem} a manutenção de \textbf{focos} e atenções simultâneas.
A \textbf{Competência} 9 está voltada ao desenvolvimento da empatia e colaboração.
Nessas dimensões, busca -se a valorização da diversidade e da alteridade, na perspectiva do acolhimento, do diálogo, da convivência e da colaboração. \textbf{Espera} -se que o estudante possa construir grupos e ambientes em diferentes contextos culturais e deles participar, combatendo o preconceito, promovendo o diálogo e o entendimento entre as pessoas em prol da melhoria da convivência em grupos culturais distintos, por \textbf{intermédio} da interação com outrem e com a diversidade.
E também que possa planejar e realizar projetos e ações com empatia e colaboração, conectando -se a grupos, comunidades e culturas locais e globais, reconhecendo as necessidades e interesses de diferentes sujeitos em contextos sociais \textbf{complexos}.
Já as dimensões da \textbf{Competência} 10 são responsabilidade, valores e cidadania, buscando o desenvolvimento de \textbf{posicionamento} consistente em relação a direitos e responsabilidades, considerando o bem coletivo, a tomada de decisão de forma colaborativa, responsável e ética, a participação social e liderança, bem como a solução de \textbf{problemas} \textbf{complexos}, pautados em princípios inclusivos, éticos, democráticos, solidários e sustentáveis.
Assim, é importante que os projetos escolares e as propostas pedagógicas ofereçam aos estudantes a oportunidade de vivenciar situações em que possam desenvolver essas habilidades, visando à formação integral e à atuação social de forma cada vez mais responsável e ética.
É importante \textbf{salientar} que todas essas \textbf{Competências} estão \textbf{inter-relacionadas} com as \textbf{competências} das áreas do conhecimento e dos componentes curriculares, que convergem para o desenvolvimento e consolidação essenciais para a formação humana e integral das \textbf{juventudes}.
Destaca -se a importância do trabalho pedagógico articulado, com intencionalidade, sistemático e progressivo visando a esses propósitos para que o estudante ao longo dos três anos do Ensino Médio aprofunde seus conhecimentos e habilidades.

% matched lemmas: abordagem, análise, aplicativo, aprimorar, argumento, articulado, caminho, capacidade, cartográfico, ciência, competência, complexo, compreensão, contemporaneidade, crítico, debate, demanda, digital, discursivo, dizer, esperar, evento, exigir, final, foco, fragilidade, imprescindível, influência, inter-relacionar, inter-relação, intercultural, intermédio, introdutório, juventude, maneira, multiletramento, ordenado, posicionamento, problema, profundidade, progressão, projeção, raciocínio, salientar, socioambiental, século, tarefa, tecnologia, verificar, vista
\end{textsample}
